\documentclass[UTF8,11pt]{ctexart}

\usepackage[a4paper,margin=2.2cm]{geometry}
\usepackage{booktabs}
\usepackage{array}
\usepackage{hyperref}
\usepackage{xcolor}
\usepackage{enumitem}

\hypersetup{
  colorlinks=true,
  linkcolor=blue,
  urlcolor=blue
}

\setlist[itemize]{leftmargin=1.6em}
\setlist[enumerate]{leftmargin=1.8em}
\renewcommand{\arraystretch}{1.18}

\title{\textbf{MMFi PA-CSI-DG Weekly Report}}
\author{mmfi-pa-csi-dg-workspace}
\date{2026-05-25 -- 2026-05-31}

\begin{document}
\maketitle

\section{本周目标}

本周目标不是继续调 PA-CSI-DG 参数，而是重新建立 MMFi 的可信 segment-level 实验协议。重点检查 sample unit、segment/window 定义、split 泄漏、duration confound、subject shift、packet/phase 表示，以及 small-model sanity gates。

\section{已完成实验操作}

\begin{enumerate}
  \item 建立独立工作区：
\begin{verbatim}
C:/Users/howie/Desktop/mmfi-pa-csi-dg-workspace/
\end{verbatim}
  原 \texttt{pa-csi-dg} 仓库只作为参考，MMFi 探索没有混入原项目。

  \item 核对 MMFi 数据源：
\begin{verbatim}
raw root:    D:/MMFi_Dataset
segment CSV: D:/datasets/mmfi/raw_drive/MMFi_action_segments.csv
\end{verbatim}
  抽查确认当前使用的 WiFi-CSI 文件和之前下载/解压的部分 MMFi 文件一致。

  \item 完成 inventory 和 segment length 检查：
\begin{center}
\begin{tabular}{cc}
\toprule
Quantile & Segment length \\
\midrule
p50 & 19 frames \\
p75 & 24 frames \\
p95 & 32 frames \\
p99 & 38 frames \\
\bottomrule
\end{tabular}
\end{center}
  因此暂时不使用 T=64/96，第一版采用 T=32。

  \item 生成第一版 segment-level feature cache：
\begin{verbatim}
feature:  P0 packet mean + A1 log1p amplitude + phase none + T32
shape:    [N, 32, 342]
segments: 16447
\end{verbatim}

  \item 生成并验证 split manifests：
\begin{verbatim}
S0 label-shuffle
S2 same-env segment-disjoint
S3 same-env subject-disjoint
S4a/S4b/S4c
S5b strict LOEO, target E01-E04
\end{verbatim}
  已检查 sample、segment、subject、target leakage。

  \item 完成 small baselines：
\begin{verbatim}
M_length duration-only
M_static static amplitude
M0 handcrafted amplitude stats
M1 TemporalCNN
\end{verbatim}

  \item 完成 A1+delta ablation：
\begin{verbatim}
feature: concat(A1_log1p_resampled, delta_time_A1_log1p_resampled)
shape:   [N, 32, 684]
\end{verbatim}

  \item 完成 weighted A1+delta TCN 的 S5b 3-seed strict LOEO：
\begin{verbatim}
seeds: 42, 43, 44
target-env mean accuracy = 0.1407 +/- 0.0186
target-env mean macro-F1 = 0.1021 +/- 0.0114
\end{verbatim}

  \item 补充 sanity gates。S0 label-shuffle 通过：
\begin{verbatim}
weighted A1+delta TCN:
  test accuracy = 0.0313
  test macro-F1 = 0.0298
\end{verbatim}
  S2/S3 显示明显 subject shift：
\begin{verbatim}
S2 E01 segment-disjoint macro-F1 = 0.9553
S3 E01 subject-disjoint macro-F1 = 0.2496

S2 E03 segment-disjoint macro-F1 = 0.8518
S3 E03 subject-disjoint macro-F1 = 0.0857
\end{verbatim}
\end{enumerate}

\section{关键结果}

\begin{center}
\begin{tabular}{>{\raggedright\arraybackslash}p{0.43\linewidth}cc}
\toprule
Experiment & Accuracy & Macro-F1 \\
\midrule
S0 label-shuffle TCN & 0.0313 & 0.0298 \\
S2 E01 segment-disjoint TCN & 0.9578 & 0.9553 \\
S3 E01 subject-disjoint TCN & 0.3202 & 0.2496 \\
S2 E03 segment-disjoint TCN & 0.8606 & 0.8518 \\
S3 E03 subject-disjoint TCN & 0.0938 & 0.0857 \\
S5b 3-seed target-env mean & 0.1407 & 0.1021 \\
\bottomrule
\end{tabular}
\end{center}

\section{主要发现}

\begin{itemize}
  \item 当前 pipeline 没有明显 label/split 泄漏。S0 label-shuffle 后 val/test 接近随机。
  \item MMFi 的 CSI/action 信号在同环境、segment-disjoint 下非常可学，说明数据、label、loader、模型训练路径不是主要问题。
  \item 一旦 subject-disjoint，性能大幅下降。subject shift 是当前最大障碍之一。
  \item S5 strict LOEO 实际是 unseen environment + unseen subject group，不是 pure environment DG。
  \item Duration confound 很强。length-only baseline 在 S5 target accuracy 上可达到约 0.17--0.22，而当前 CSI TCN accuracy 约 0.14。
  \item A1+delta 明显优于原始 P0/A1，但仍不足以支撑 PA-CSI-DG 或 SupCon 结论。
\end{itemize}

\section{当前问题}

\begin{enumerate}
  \item Strict LOEO 表现很差：
\begin{verbatim}
best small-model S5b macro-F1 mean = 0.1021
best small-model S5b accuracy mean = 0.1407
\end{verbatim}
  这只是略高于随机，远不是可接受的最终结果。

  \item Subject shift 严重。S2 很高但 S3 急剧下降，说明模型学到大量 subject-specific pattern。

  \item Duration / segment length confound 还没有控制住。

  \item 当前 delta 是 fixed-T 派生 delta，不是 raw segment 上的 converter-level delta。

  \item Packet 维信息还没有充分利用，目前主结果仍基于 P0 packet mean。

  \item Phase V1/V2 还没进入可信实验，phase-only gate 未完成。

  \item 还没有进入 full PA-CSI-DG 或 SupCon 的条件。
\end{enumerate}

\section{下周建议}

优先继续 protocol 和 preprocessing，不建议直接上 full PA-CSI-DG。建议顺序：

\begin{enumerate}
  \item 做 duration/length-controlled evaluation。
  \item 实现 converter-level raw A1\_delta。
  \item 跑 P1 packet mean/std/diff 和 P3a packet dynamics。
  \item 做 per-class confusion / failure analysis。
  \item 做 subject-ID predictability diagnostic。
  \item 再决定是否进入 phase V1 和 full PA-CSI-DG。
\end{enumerate}

\section{结论}

\begin{quote}
Pipeline 可用，但 MMFi strict DG 任务尚未解决。当前失败点主要是 subject shift、duration confound 和 P0 feature 表示不足。
\end{quote}

\end{document}
